\section{Discussion} \subsection{Interpretation of the main wavelet findings} This study extends the interpretation of urban park air-quality effects beyond static inside--outside mean concentration differences by showing that Bangkok parks modify PM$_{2.5}$ variability in a strongly timescale-dependent manner. The wavelet attenuation--coherence framework reveals that parks do not simply reduce PM$_{2.5}$ by a fixed amount. Rather, they selectively damp particular temporal rhythms while still preserving part of the broader urban signal. In the present 2021 case study, damping was strongest and most consistent from sub-diurnal to multi-day scales, but the band structure was not monotonic: both the 2--6 h and diurnal bands were weaker, and the diurnal band showed the least damping overall. At the shortest 2--6 h band, a small number of parks even showed positive attenuation, indicating local amplification rather than damping at that timescale. This distinction matters scientifically and practically. A park that suppresses multi-day episode structure or sub-diurnal fluctuation amplitude is functioning differently from one that only marginally alters fast fluctuations or the daily cycle, even if their mean PM$_{2.5}$ values appear similar over time. The central methodological contribution of this study is therefore not just the use of continuous wavelet transform (CWT), but the pairing of attenuation and coherence as complementary descriptors of park buffering behavior. The regime concept is the key substantive novelty. The eight retained 2021 park--outside pairs did not behave as a single homogeneous class; instead, they organized naturally into strong decouplers, moderate broadband dampers, and weak followers or mixed cases. This shifts the framing of the urban-park question. Asking whether parks help with air quality is too blunt. The more useful question is which parks damp which temporal rhythms, how strongly they remain coupled to surrounding urban dynamics, and under what spatial conditions those behaviors emerge. In this sense, the regime classification is more informative than a citywide average effect size, because it respects internal heterogeneity among parks. It also provides a practical interpretation bridge between the spectral results and urban design: different park types may contribute differently to the buffering of short-term fluctuations, daily rhythms, or broader episode-scale variability. A stricter 14-day boundary-trim sensitivity check did not materially alter this qualitative three-regime structure, although some long-band attenuation magnitudes changed moderately. A second implication of the wavelet results is that amplitude damping and temporal coupling are not interchangeable concepts. From an exposure perspective, this means that a park may reduce the intensity of PM$_{2.5}$ episodes experienced inside the park even when the timing of broader urban pollution rhythms remains largely shared with the surrounding city. In practical terms, this kind of buffering may still be highly meaningful for public health and park use. Lowering the amplitude of daily or multi-day pollution fluctuations can reduce the severity of peak exposure conditions for park users even if the park does not become fully disconnected from city-scale atmospheric behavior. Thus, the value of a park as an air-quality buffer does not depend on becoming an isolated micro-atmosphere; it may instead lie in moderating the strength of pollution variability while remaining temporally embedded within the larger urban system. Two parks are especially revealing because they resist simplistic explanation. Nong Chok Park behaved as a strong decoupler despite not being reducible to a single mean greenness signal. This suggests that favorable spatial configuration, radial contrast structure, or local source separation may matter more than average vegetation alone. At the other end, Lumpini Park did not emerge as one of the strongest dampers despite its prominent green identity. This is a useful reminder that visual or intuitive greenness does not guarantee strong spectral damping performance. These cases strengthen, rather than weaken, the argument for a multi-factorial explanatory framework. They show why a one-variable interpretation would have been misleading and why the layered combination of mean land-cover contrast, ring gradients, water-related context, and morphology is more appropriate. \subsection{Spatial explanatory hierarchy} A major result of this study is that average greenness alone is not an adequate explanatory metric. The simple buffer-based results did show meaningful first-order relationships for park--outside NDVI and NDBI contrast, but these were not sufficient to explain the full hierarchy of damping behavior. The annulus-based analysis strengthened the interpretation considerably by showing that radial structure mattered more than mean composition alone. This is physically plausible. PM$_{2.5}$ attenuation is not expected to depend simply on whether vegetation is present, but on how it is organized: whether a protected vegetated interior exists, how quickly built-up pressure increases toward the edge, whether contrast persists outward from the park core, and how those transitions compare with the outside reference site. In that sense, the concepts of green-core advantage, built-edge pressure, and contrast persistence are not merely statistical labels; they are spatially intuitive descriptors of how a park is embedded within the surrounding urban fabric. Because NDVI and NDBI were extracted from station-centered buffers rather than clipped strictly to park boundaries, the built-up signal observed near park stations may reflect either internal hardscape or proximity to surrounding urban land cover, and these two contributions cannot be fully disentangled in the present dataset. The importance of ring-gradient structure also helps explain why the strongest explanatory signals often appeared for coherence, especially at short and sub-diurnal timescales. Short-band coherence reflects how tightly the park and outside station move together in time, even if their amplitudes differ. The finding that outside-site radial land-cover structure was strongly associated with short-band coherence suggests that local comparator context strongly influences whether a park remains dynamically synchronized with nearby urban fluctuations. This is a notable result because it shows that park performance is relational. A park is not acting in isolation; its apparent damping strength depends partly on what it is being compared against. If the outside site has a greener core, a sharper built-up edge, or a different water-related spatial organization, the inside--outside comparison itself changes in meaning. That outside-comparator insight is one of the most interesting findings of the study. Much of the literature implicitly treats the outside site as a neutral baseline, but the present results suggest that the baseline itself carries spatial structure that shapes the observed inside--outside relationship. This helps explain why some parks may appear to perform differently depending on their comparator context. It also means that future paired park studies should treat comparator selection as a substantive part of study design rather than as a purely technical detail. In the present workflow, the pairing was based on nearest geodesic distance under a transparent quality gate, which supports reproducibility. However, the results show that the spatial context of the outside site is not passive; it is an active part of the explanatory system. The role of water and morphology should be interpreted within that layered framework. Water-related metrics added a clear and targeted coherence-oriented signal, especially for short-band and 6--18 h synchrony, suggesting that moisture-sensitive local context may affect how quickly parks and outside sites respond together to changing conditions. However, water did not replace land-cover gradients as the dominant explanatory layer for attenuation across all bands. Morphology, even after reviewed polygon correction, remained secondary. This is an important negative result. It suggests that park footprint shape alone is not the primary control on PM$_{2.5}$ buffering behavior across this sample. In addition, the negative associations involving park perimeter should not be interpreted as evidence that greater edge exposure is beneficial. In this small sample, perimeter likely co-varies with park area and the existence of a larger protected vegetated core, so it is more plausibly acting as a partial size proxy than as a pure edge metric. Park form therefore becomes most meaningful when interpreted together with surrounding land-cover structure and comparator-site context. In other words, geometry matters, but not in isolation. Taken together, the explanatory hierarchy was clear. Mean land-cover contrast provided a useful first-order signal; annulus-based land-cover gradients provided the broadest explanatory improvement across attenuation and coherence; water-related context contributed some of the strongest individual associations, but in a narrower and more coherence-focused domain; and morphology offered secondary structural support. No single predictor explained all eight parks completely. Rather than a weakness, this is the result itself: PM$_{2.5}$ buffering by urban parks in Bangkok depends on a layered combination of internal green structure, surrounding built-up pressure, comparator-site context, water-related microenvironment, and park form. This layered explanatory hierarchy is synthesized conceptually in Fig.~\ref{fig:hierarchy_pyramid}. \begin{figure}[htbp]
\centering
\resizebox{\textwidth}{!}{%
\begin{tikzpicture}[font=\sffamily]

% =========================
% Color Definitions
% =========================
\colorlet{col1}{blue!18}
\colorlet{col2}{green!18}
\colorlet{col3}{teal!22}
\colorlet{col4}{gray!20}

\colorlet{line1}{blue!60!black}
\colorlet{line2}{green!50!black}
\colorlet{line3}{teal!70!black}
\colorlet{line4}{gray!70!black}

% =========================
% Pyramid Geometry
% =========================
\def\W{9.5}   % Slightly widened base to easily fit the \normalsize text
\def\H{9.6}   % Total height
\def\CX{-3.0} % Center X of pyramid to balance the boxes on the right

\coordinate (BL) at (\CX-\W/2, 0);
\coordinate (BR) at (\CX+\W/2, 0);
\coordinate (TOP) at (\CX, \H);

% Y-coordinates of slice boundaries
\def\yA{2.4}
\def\yB{4.8}
\def\yC{7.2}

% Calculate coordinates on the edges
\coordinate (L1) at ($(BL)!\yA/\H!(TOP)$);
\coordinate (R1) at ($(BR)!\yA/\H!(TOP)$);
\coordinate (L2) at ($(BL)!\yB/\H!(TOP)$);
\coordinate (R2) at ($(BR)!\yB/\H!(TOP)$);
\coordinate (L3) at ($(BL)!\yC/\H!(TOP)$);
\coordinate (R3) at ($(BR)!\yC/\H!(TOP)$);

% =========================
% Filled layers & Borders
% =========================
\fill[col1] (BL) -- (BR) -- (R1) -- (L1) -- cycle;
\fill[col2] (L1) -- (R1) -- (R2) -- (L2) -- cycle;
\fill[col3] (L2) -- (R2) -- (R3) -- (L3) -- cycle;
\fill[col4] (L3) -- (R3) -- (TOP) -- cycle;

% Internal layer borders
\draw[white, line width=1.5pt] (L1) -- (R1);
\draw[white, line width=1.5pt] (L2) -- (R2);
\draw[white, line width=1.5pt] (L3) -- (R3);

% Outer pyramid border
\draw[black!60, line width=1pt] (BL) -- (BR) -- (TOP) -- cycle;

% =========================
% Box Styles
% =========================
\tikzset{
    infobox/.style={
        draw,
        rounded corners=2.5pt,
        line width=0.8pt,
        fill=white,
        align=left,
        text width=4.5cm,
        inner sep=6pt,
        font=\footnotesize % Reduced to standard diagram font size
    }
}

% =========================
% Edge Connectors & Boxes
% =========================
\coordinate (M1) at ($(BR)!0.5!(R1)$);
\coordinate (M2) at ($(R1)!0.5!(R2)$);
\coordinate (M3) at ($(R2)!0.5!(R3)$);
\coordinate (M4) at ($(R3)!0.5!(TOP)$);

\def\boxX{2.2} % Left edge of all boxes

% Box 1 (Mean)
\node[infobox, draw=line1] (box1) at (\boxX, 1.2) [anchor=west] {
    \textbf{Target domain:} attenuation\\[0.1cm]
    \textbf{Relative role:} useful first-order signal
};
\draw[line1, line width=0.9pt] (M1) -- (box1.west);

% Box 2 (Ring)
\node[infobox, draw=line2] (box2) at (\boxX, 3.6) [anchor=west] {
    \textbf{Target domain:} attenuation + coherence\\[0.1cm]
    \textbf{Relative role:} broadest explanatory improvement
};
\draw[line2, line width=0.9pt] (M2) -- (box2.west);

% Box 3 (Water)
\node[infobox, draw=line3] (box3) at (\boxX, 6.0) [anchor=west] {
    \textbf{Target domain:} coherence-focused\\[0.1cm]
    \textbf{Relative role:} strong but narrower contribution
};
\draw[line3, line width=0.9pt] (M3) -- (box3.west);

% Box 4 (Morphology)
\node[infobox, draw=line4] (box4) at (\boxX, 8.4) [anchor=west] {
    \textbf{Target domain:} structural support\\[0.1cm]
    \textbf{Relative role:} secondary explanatory layer
};
\draw[line4, line width=0.9pt] (M4) -- (box4.west);

% =========================
% Internal Pyramid Labels
% =========================
\node[font=\normalsize\bfseries] at (\CX, 1.2) {Mean land-cover contrast};
\node[font=\normalsize\bfseries] at (\CX, 3.6) {Ring-gradient structure};
\node[font=\normalsize\bfseries] at (\CX, 6.0) {Water context};
% Morphology is given a smaller font and lowered slightly to safely clear the borders
\node[font=\scriptsize\bfseries] at (\CX, 8.1) {Morphology};

% =========================
% Grouping Brace
% =========================
\draw[decorate,decoration={brace,amplitude=5pt,mirror},line width=1pt,black!70]
(7.0, 0.4) -- (7.0, 4.4)
node[midway, right=10pt, align=left, font=\small\bfseries, text width=2.8cm, color=black!85]
{Dominant\\[2pt]land-cover\\[2pt]explanatory\\[2pt]base};

\end{tikzpicture}%
}
\caption{Conceptual synthesis of the layered explanatory hierarchy identified in the Bangkok park--outside PM$_{2.5}$ wavelet analysis. Mean land-cover contrast provided a useful first-order signal, ring-gradient structure provided the broadest explanatory improvement across attenuation and coherence, water context contributed strong but narrower coherence-focused associations, and morphology acted as a secondary structural layer rather than a dominant control.}
\label{fig:hierarchy_pyramid}
\end{figure} \subsection{Implications for urban park design and management} Although the present analysis is exploratory and should not be interpreted as a direct intervention trial, the results do suggest several practical implications for urban park design and management in Bangkok. The first is that total green area alone is unlikely to be an adequate design target. The explanatory hierarchy developed here indicates that spatial organization matters at least as much as gross park area. In practical terms, parks with a more protected vegetated interior and lower central built-up pressure are likely to provide more favorable conditions for PM$_{2.5}$ damping than parks whose area is dominated by fragmented or edge-exposed green cover. A second implication is that the center of a park should be treated as a protected buffering zone rather than as the preferred location for extensive hardscape. The annulus-based results suggest that when built-up pressure penetrates toward the interior, the park loses part of its damping advantage. This does not mean that paved walkways, event spaces, or service infrastructure should be eliminated, but it does imply that large concrete plazas, wide paved corridors, or extensive impervious surfaces placed deep within the central zone may erode the function of the park as a PM$_{2.5}$ buffer. A more effective layout would place hardscape more selectively and preserve a comparatively uninterrupted interior vegetation core. A third implication concerns the park edge. The ring-gradient analysis indicates that transitions between the surrounding built environment and the park interior are highly relevant to buffering behavior. This suggests that the outer park zone should not be treated as a simple decorative boundary. Instead, boundary planting and edge design may be important for reducing direct penetration of nearby urban fluctuations into the interior. In practical planning terms, dense and vertically structured edge vegetation may be more valuable than sparse tree lines if the objective is to strengthen the contrast between the park interior and the surrounding street environment. A fourth implication concerns water features. In the present study, water-related metrics were most strongly associated with short-band and 6--18 h coherence rather than with universal broadband attenuation. This suggests that ponds, canals, or moisture-related landscape elements should not be treated as direct substitutes for vegetation, nor as simple particulate filters. However, the results are consistent with the possibility that water-sensitive local context helps stabilize short-timescale microenvironmental behavior and modifies the degree to which park and outside sites fluctuate together. In design terms, well-maintained water bodies may therefore contribute supportive microclimatic value when integrated with vegetation and favorable land-cover structure. A fifth implication is that park performance cannot be separated from the surrounding urban source environment. The relational importance of the outside comparator implies that heavy traffic, transport bottlenecks, or intense built-up pressure immediately adjacent to the park perimeter may overwhelm part of the park's buffering advantage. This is particularly relevant for cases such as Lumpini Park, where strong park identity does not automatically translate into the strongest damping regime. The practical lesson is that park design and traffic management should not be treated as independent policy domains. A park may perform better when perimeter traffic pressure, near-edge source density, and surrounding hardscape intensity are considered as part of the same environmental design problem. Taken together, these implications suggest a shift in planning emphasis: from maximizing green area in the abstract toward protecting vegetated cores, reducing central hardscape pressure, strengthening park-edge transitions, integrating supportive water context, and managing the immediate urban perimeter. These recommendations should be treated as design implications consistent with the present data, not as definitive prescriptions. However, they show how a wavelet-based, timescale-specific framework can be translated into concrete urban-management thinking more effectively than a static mean-concentration comparison alone. \subsection{Limitations and future directions} Several limitations should be stated explicitly. First, the effective sample size was small ($n = 8$ retained pairs), so all statistical relationships must be treated as exploratory. The Spearman and leave-one-out results are valuable as robustness-oriented pattern tests, but they do not support strong causal inference. Second, the paper is intentionally framed as a 2021 Bangkok case study. This is not because later years were ignored, but because 2022--2023 failed the same defensible annual quality gate, largely due to insufficient hourly completeness in the outside-station records. That decision strengthens the credibility of the paper, but it also means the conclusions should be interpreted as rigorously quality-gated for 2021 rather than as a multi-year climatology. Third, three parks shared the same outside comparator (Bang Sue District), which may have increased apparent similarity among some regime members and means that those parks should not be interpreted as fully independent confirmations of the same regime pattern. Fourth, the wavelet summaries compress behavior across a full year and therefore cannot resolve every local source event, directional transport pathway, or shifting meteorological regime within that window. Finally, the land-cover, water, and morphology features are correlative rather than causal. They identify spatial associations with damping behavior, but they do not, by themselves, prove mechanism. In addition, the main wavelet summaries were based on a fixed 24-h boundary trim rather than full scale-dependent cone-of-influence masking, so long-period attenuation estimates should be interpreted more cautiously than short-period or coherence-based results; a sensitivity check using a stricter 14-day trim changed some long-band attenuation magnitudes moderately, but the qualitative regime structure and the broader interpretation remained broadly consistent. These limitations point directly to future directions. The most immediate need is replication using denser and more stable monitoring networks with higher-quality outside comparators across multiple years. A second priority is the inclusion of directional or anisotropic spatial features, such as upwind--downwind sectors, because radial averages may miss important transport geometry. A third extension would be tighter integration with meteorological variables, especially wind speed, wind direction, atmospheric stability, and boundary-layer structure, which are likely to modulate when and how park damping becomes visible. A fourth future direction is application to other cities, where the same attenuation--coherence framework could test whether the Bangkok regime structure is generalizable or whether it depends on local urban morphology and climate. Despite these limitations, the broader implication is clear. Urban parks in Bangkok did not behave as interchangeable green spaces; they buffered PM$_{2.5}$ variability in different ways across timescales, and those differences were linked more strongly to layered spatial structure than to average greenness alone. The study therefore suggests that policy and design should move beyond the simple assumption that more green area automatically produces stronger air-quality benefits. Protected interior vegetation, reduced built-up edge pressure, and favorable spatial integration with surrounding urban land cover are likely to matter more than area alone. More broadly, the wavelet attenuation--coherence framework presented here offers a reproducible and transferable way to evaluate how urban green spaces interact with air-pollution dynamics in time, not just in mean concentration space.